\documentclass[10pt]{ctexart}
\usepackage[a4paper,margin=1in]{geometry}
\usepackage{amsmath,amssymb,booktabs,longtable}
\usepackage{enumitem}
\usepackage{microtype}
\usepackage{hyperref}
\setlist[itemize]{leftmargin=1.5em}
\setlist[enumerate]{leftmargin=1.8em}
\title{001：StepPPO 研究聊天记录归档}
\author{用户与 Codex}
\date{2026-06-05}
\begin{document}
\maketitle
\vspace{-1.2em}

\section*{说明}
本文归档当前会话中围绕 Step-wise PPO、GRPO、GiGPO、实验脚本、可视化、文档和算法改进的主要讨论内容。为避免 PDF 过长，本文保留用户问题、关键回答、形成的决策、代码与文档产物；对长日志和工具输出只保留结论与路径。

\section*{一、研究目标与初始规划}
用户首先提出希望做一个 step-wise PPO 实验，并要求先理解当前 PPO、GRPO 和 GiGPO 的实现。经过对齐后，形成如下目标：
\begin{itemize}
\item 任务是 agentic multi-step episode，每个 episode 包含多个 environment steps。
\item 每个 step 中，agent 接收 context/observation，生成一个完整 action，例如 tool call 或文本动作。
\item Step-wise PPO 的 value、return、TD error 和 GAE 都在 environment-step 维度计算，而不是 token 维度。
\item 每个 action 只产生一个 value prediction；底层模型可以每个 token 输出 value，但主实现只取 action 生成前边界处的 scalar value。
\end{itemize}

\section*{二、Critic 与 Value Head 的共识}
关于 PPO critic，讨论中形成的共识是：
\begin{itemize}
\item 传统 PPO 可以维护独立 critic 模型，但本实验不希望额外维护一套 critic backbone。
\item 当前 StepPPO 采用 shared-backbone actor-critic：在 actor CausalLM 上额外挂一个 TRL-style value head。
\item value head 是小型 regression head，主版本为 MLP，输入为 actor hidden state，输出 scalar value。
\item 默认 value position 是 \path{pre_action_last_context_token}，即 action 生成前最后一个 context token，用来估计严格的 $V(s_t)$。
\item 取 first response token 或 last response token 更接近 action-conditioned baseline，保留为未来 ablation，不作为主定义。
\end{itemize}

严格 step value 形式为：
\[
V_{\theta,\psi}(s_t)=v_\psi(H_t[b_t]),
\qquad b_t=n_t-L_t-1,
\]
其中 $n_t$ 是 context+response 总长度，$L_t$ 是 response 长度。

\section*{三、StepPPO Advantage 的定义}
用户强调 StepPPO 的 return、GAE、TD-error 都应沿 step 维度计算。最终采用：
\[
\delta_t=r_t+\gamma(1-d_t)V^{old}(s_{t+1})-V^{old}(s_t),
\]
\[
A_t^{step}=\delta_t+\gamma\lambda(1-d_t)A_{t+1}^{step},
\qquad
R_t^{step}=A_t^{step}+V^{old}(s_t).
\]

同时保留 episode-level GRPO signal：
\[
A_i^{epi}=\frac{R_i^{epi}-\mu_g}{\sigma_g+\epsilon}.
\]

v1 版本定义为：
\[
A_t^{final,v1}=\frac{\widehat A_t^{epi}+w\widehat A_t^{step}}{1+w}.
\]

v0 base 版本定义为：
\[
A_t^{final,v0}=\frac{\widehat A_t^{epi}+wA_t^{step}}{1+w}.
\]

\section*{四、Baseline 选择讨论}
围绕 baseline，形成如下结论：
\begin{itemize}
\item GRPO 是 non-step baseline，因为它只使用 episode-level grouped relative advantage。
\item GiGPO 是更重要的 step-aware baseline，因为它已经利用 step return、anchor observation grouping 和 step-level relative advantage。
\item StepPPO 应主要和 GiGPO 对比，回答 learned step-level critic correction 是否优于 critic-free step grouping。
\item 如果想衡量 step 信息本身是否有用，可以对比 GRPO。
\end{itemize}

\section*{五、v1、v0 与实验顺序}
用户重新规划实验后，明确：
\begin{itemize}
\item GiGPO 三个 seed 已经跑过。
\item GRPO 需要作为额外 baseline 补跑。
\item StepPPO 准备两个版本：v1\_step\_norm 与 v0\_base。
\item 首轮 seed 使用 2026，顺序为 StepPPO-v1、StepPPO-v0、GRPO。
\end{itemize}

对应批量脚本为：
\begin{verbatim}
exps/run_webshop_step_ppo_v1_v0_grpo_seed2026.sh
\end{verbatim}

\section*{六、可视化与已有训练结果}
用户要求将 StepPPO 与 GiGPO 三个 seed 的结果可视化对比。生成的主要输出包括：
\begin{itemize}
\item \path{VISULIZATION/figures/fig_gigpo_vs_step_ppo_val_metrics.png}
\item \path{VISULIZATION/figures/fig_step_ppo_instability_diagnostics.png}
\item \path{VISULIZATION/figures/fig_step_ppo_time_resource_diagnostics.png}
\end{itemize}

已有诊断显示：
\begin{itemize}
\item StepPPO-v1 相对 GiGPO validation 指标更低。
\item StepPPO-v1 valid action ratio 更低。
\item StepPPO-v1 response clipping 更高。
\item StepPPO-v1 value loss 和 grad norm 更大。
\item StepPPO-v1 的 rollout generation 更慢，后期非验证 step 约为 GiGPO 的 $2.02\times$。
\end{itemize}

\section*{七、StepPPO 训练慢的定位}
对时间和资源消耗的定位结论：
\begin{itemize}
\item 主因不是显存不足，StepPPO 平均 GPU allocated memory 反而低于 GiGPO。
\item 最大时间差来自 rollout generation：StepPPO 生成更长、更容易截断、更可能 invalid 的 response。
\item 直接实现开销也存在：StepPPO 需要 value prediction 和 value loss。
\item 原 old-log-probability 阶段存在 logprob 与 value 分两次 forward 的开销。
\end{itemize}

时间诊断文档为：
\begin{verbatim}
EXPS/010_step_ppo_time_resource_analysis.pdf
\end{verbatim}

\section*{八、文档体系整理}
用户要求 EXPS 下文档以 TeX/PDF 形式保存，并指定除 method/算法介绍用英文外，其它以中文为主。当前约定为：
\begin{itemize}
\item \path{EXPS/tex/} 存放实验文档 TeX 源文件。
\item \path{EXPS/} 根目录存放对应 PDF。
\item \path{003_shared_value_head_step_ppo_method} 保持英文 method 文档。
\item 其它实验规划、实现说明、诊断分析文档以中文为主。
\end{itemize}

已提交过的相关 commit 包括：
\begin{itemize}
\item \texttt{ede5d07 Add TeX experiment docs and log conventions}
\item \texttt{8279bea Localize experiment docs except method}
\end{itemize}

\section*{九、StepPPO-v2 改进}
用户询问是否可以改进 StepPPO 后，做了一个低风险 v2 改进版本：
\begin{enumerate}
\item old-log-probability 阶段改为 single forward 同时返回 logprob 和 state value，减少重复计算。
\item 新增 \path{normalize_final_advantage} 配置，默认关闭。
\item 新增 v2 stable 脚本，使用 detached value backbone 和更小 value loss coefficient。
\end{enumerate}

v2 stable 核心配置为：
\begin{verbatim}
algorithm.step_ppo.normalize_step_advantage=True
algorithm.step_ppo.normalize_final_advantage=True
actor_rollout_ref.actor.value_head.detach_value_backbone=True
actor_rollout_ref.actor.value_head.value_loss_coef=0.03
\end{verbatim}

新增脚本为：
\begin{verbatim}
exps/run_webshop_step_ppo_v2_stable_4gpu_paper_align_simple.sh
\end{verbatim}

提交为：
\begin{verbatim}
c5829ff Improve StepPPO stability and old value collection
\end{verbatim}

\section*{十、最近一次算法细节讨论}
用户最后询问两个问题：
\begin{enumerate}
\item 平衡 $A_t^{epi}$ 与 $A_t^{step}$ 的 $w$ 默认是多少？
\item 当前 v1 除了加入 $A^{step}$ 和 value loss 外，是否与 GiGPO 完全对齐？
\end{enumerate}

回答结论：
\begin{itemize}
\item $w$ 默认是 $1.0$，即 StepPPO-v1 中 episode:step 的比例是 $1:1$。
\item 但 StepPPO-v1 用 $1+w$ 做平均，而 GiGPO 当前是直接相加：$A^{GiGPO}=A^{epi}+wA^{step}$。
\item StepPPO-v1 并不与 GiGPO 完全对齐。主要差异包括 episode normalization、step advantage 来源、step normalization 粒度、final advantage scale 和 microbatch 配置。
\end{itemize}

更准确的描述是：StepPPO-v1 是以 GRPO-style episode advantage 为 anchor，再加入 critic/GAE step correction 的 shared-backbone actor-critic 方法；它借鉴了 GiGPO 的 episode+step 框架，但不等价于 GiGPO 加一个 value loss。

\section*{十一、后续建议}
如果希望做更严格的 GiGPO-aligned StepPPO，建议开新版本，例如 \texttt{v3\_gigpo\_aligned}：
\begin{itemize}
\item episode advantage 语义对齐 GiGPO 的 \texttt{mean\_norm} 或 \texttt{mean\_std\_norm}。
\item final advantage 与 GiGPO 一样使用 $A^{epi}+wA^{step}$，或者显式做 scale ablation。
\item step normalization 不做全 batch whitening，而考虑按 \path{anchor_obs} group 或至少按 \path{uid} group。
\item 只把 critic/GAE step signal 作为变量，其它尽量固定。
\end{itemize}

\end{document}
